\chapter[African Cities, States, and Knowledge]{What Cities, States, and Knowledge Traditions Did Africa Possess?}
\section{A Manuscript Rescued in a Rice Sack}
First pick up an old manuscript page.

About A4 size, it has wormholes at its edges and paper yellowed like tea stains. Arabic letters cover it, with occasional red lines among black ink like a teacher's emphasis. Originally it was no bound book: in West Africa centuries ago, paper cost more than books, and learning often occupied loose sheets gathered with leather ties or cloth covers and passed down generations.

The page comes from Timbuktu in today's Mali. Its library had no building, desk, or formal name, only a wooden chest beneath a family's wardrobe, holding hundreds of manuscripts on mathematics, astronomy, medicine, law, poetry, accounts, and correspondence. In 2012, armed groups occupied the city and destroyed tombs and holy sites. Widely reported accounts describe librarians and hundreds of families risking their lives to smuggle roughly 350,000 manuscripts south in rice and flour sacks, using donkey carts and canoes. Your page may have jolted inside such a sack.

Remember three facts. It comes from the Sahara's southern edge. It documents substantial scholarly life. And rescuers still risked their lives for it in 2012, because the history it represents had long been declared nonexistent.

Our question lies in the paper: what cities, states, and intellectual traditions existed in precolonial Africa, particularly south of the Sahara? Why can many teenagers name several old European cities but no African one?

\section{A Caravan Enters the Gold-and-Salt Market}
Now enter a scene.

Time: a mid-fourteenth-century morning. Place: sandy ground outside Timbuktu. The Niger makes a great bend not far north of town, with endless sand beyond its green banks.

Camel bells arrive first. Hundreds of animals file into the city after almost two months crossing the Sahara from the north. Each bears two translucent, pink-veined salt slabs from Taghaza, deep in the desert. There salt forms the ground, walls, and houses themselves.

Near the gate, camels kneel, drivers brush sand from robes, and salt goes to the scales. The rule is startlingly simple: salt is almost worth its weight in gold, no exaggeration. West African grasslands and forest margins have gold but lack salt, essential in heat and for preserving meat. Desert salt and forest gold dust meet in river markets like Timbuktu.

Follow a salt merchant and see more movement: southern kola nuts for stimulation, ivory, leather, and enslaved people; northern cloth, horses, dates, and copperware. Payments vary: gold dust weighed from leather pouches, small purchases in cowries imported from the Indian Ocean. Four or five languages bargain simultaneously, while administrators, accountants, and contract writers use Arabic.

Beyond the market lies a quieter quarter. A young man sits cross-legged in mosque shade with a writing board on his knees, copying law. His teacher requires finishing one borrowed volume before taking another. Books are scarce; local saying makes them dearer than salt or gold. Dozens of teachers live around Sankore and other great mosques, drawing students from across sub-Saharan Africa, like Europe's later students travelling to Paris and Bologna.

Caravan, market, and copyists belonged to a real world contemporary with China's late Yuan and early Ming. Soon we will read an eyewitness. First address a difficulty: this world was long declared never to have existed.

\section{A Question Already Supposedly Answered}
State the conflict before conclusions.

Compare the African and European page counts in an ordinary popular world history. The imbalance may startle you. For centuries, claims that sub-Saharan Africa had no history circulated globally: no cities, states, writing, or noteworthy events, only tribes, jungle, and unchanging primitive life.

One famous spokesman was nineteenth-century German philosopher Hegel. His philosophy-of-history lectures excluded Africa from world history, broadly arguing that it lacked states and reason. Notice the logic: having read scarcely any African primary sources, he could still conclude this because absence of history was his premise, not a research finding.

Even in a 1960s broadcast, a prominent Oxford historian said Africa might eventually offer history to teach, but currently contained only Europeans' history there; the rest was darkness, not history's subject. By then Ghana, Mali, and Songhay chronicles had been published and Great Zimbabwe's archaeological reports available for decades. Evidence existed; some refused it recognition.

Our real question is not whether African cities, states, and knowledge existed. Evidence has long answered yes, abundantly. Instead ask:

\textbf{First, why were they denied recognition as history for so long?}

\textbf{Second, if no writing means no history is a false standard, how can buried histories be recovered?}

The second matters beyond Africa. Most people in most places for most of time left no written records. If writing alone grants admission, much of humanity's past becomes officially nonexistent.

\section{Four Voices Narrating Africa}
To answer the first question, hear different parties. At least four voices describe the same continent.

\textbf{Voice one: colonisers---nothing existed before we arrived.}

Its function is obvious: occupying ahistorical land becomes bringing history rather than taking someone else's. European powers met in Berlin in 1884--1885 and partitioned Africa with rulers on maps, drawing most modern borders without any African present. The territorial drawing needed a matching story, so Africa became a blank sheet.

An extreme episode concerns Great Zimbabwe. This enormous dry-stone granite complex in today's Zimbabwe has walls over ten metres high enclosing royal and ritual spaces and once housed over ten thousand people. Late-nineteenth-century white colonisers refused African authorship, preferring the Queen of Sheba, Phoenicians, or a lost white tribe. Scientific excavations in 1905 and 1929 clearly identified medieval African dates and workmanship. Colonial officials disliked implications that Africans could build cities and colonisation was dispossession. For decades, scholars holding the correct conclusion were marginalised while textbooks retained Sheba's legend.

\textbf{Voice two: oral custodians---we are the history books.}

West African griots, sometimes called bard-historians, memorise royal descent, riverside battles, and family feuds from childhood. The Sundiata epic tells how a thirteenth-century prince defeated enemies and founded Mali. For centuries, successive griots transmitted it without paper, until twentieth-century recording and transcription. They consciously describe themselves as kings' memory: without them, ancestors' names die.

\textbf{Voice three: foreign travellers in Arabic and Chinese.}

From the eighth century, Arabic geographers, merchants, and judges wrote about West Africa. Eleventh-century Andalusian al-Bakri never left al-Andalus but interviewed traders and envoys to describe Ghana's palace, ministers' clothes, markets, and ownership of gold nuggets. Fourteenth-century Moroccan Ibn Battuta travelled Mali and the East African coast and recorded them. Chinese sources also exist: Zheng He's fleets reached East Africa, and accompanying Ma Huan's Overall Survey of the Ocean's Shores described city-states including Mogadishu in today's Somalia.

\textbf{Voice four: the earth---archaeology.}

Walls, pottery, foundations, mines, burials, imported ceramics, and beads cannot speak, but neither change their account to flatter someone. Chinese celadon and Persian pottery at Great Zimbabwe connect the inland city through coastal ports to Indian Ocean trade. Beneath Niger inland-delta mounds lies Jenne-jeno, growing into a city two thousand years ago. Layers of streets and burials quietly establish urban building before outsiders recorded it.

These voices contradict and complement one another. Practise strengthening the apparently weakest: oral tradition. Why count speech as evidence? Memory errs, stories exaggerate, genealogies flatter current kings, and centuries of transmission change tales. Every concern is valid. Griots adapt praise to patrons, and Sundiata's enemy appears almost a magical demon king.

But the same doubts apply to writing. Victors compose chronicles, travellers repeat hearsay, archives serve power. Applying equal scepticism to spoken and written sources is method; doubting only speech is prejudice. That distinction leads to our tool.

\section[Thinking Bridge: Cross-Check Evidence]{Thinking Bridge: Let Three Kinds of Evidence Check One Another}
Take away the tool of \textbf{corroboration}.

For Africa and other sparsely documented places, historians often have oral traditions, archaeology, and outsiders' writing. Each has strengths and blind spots, like three visually impaired observers around an elephant. Do not simply choose one; cross-check them with four questions:

\begin{itemize}
\item \textbf{Who produced it, from what position?} Griots serve courts, travellers are outsiders, pottery is neutral.
\item \textbf{How much later than the event is it?} Immediate records, inherited memories, and three-hundred-year-later epics carry different weight.
\item \textbf{Whom does it flatter?} Kings, readers at home, or colonial governments?
\item \textbf{What other evidence could confirm or refute it?} If an epic describes a conquest, can archaeology find contemporary walls there?
\end{itemize}
Apply it to a real case.

Mali's foundation: Sundiata's oral epic places his victory and founding in the early thirteenth century. Arabic chronicles, including later local works such as the Tarikh al-Sudan, likewise record founders' descent and campaigns within a matching chronology. Archaeologists searching the upper Niger for legendary Niani found medieval settlement and trading remains. Its capital status remains disputed, but flourishing large settlements in this region during the thirteenth to fifteenth centuries are not. Time, lineage, and broad location agree sufficiently to establish Mali's foundation historically, while a flying demon king belongs to literature.

Contradictions also supply clues. Traditional accounts attribute Ghana's decline to eleventh-century conquest by desert holy warriors. Some modern scholars question the conquest's scale and emphasise long-term drought and rerouted trade instead. Debate continues. The point is that oral memory favours \textbf{dramatic events}, heroes and wars, while slow causes often need other evidence. Each source has its own appetite for memory, making corroboration essential.

Silent trade also left evidence. Herodotus's Histories, Book IV, describes Carthaginians putting goods on the Libyan shore and withdrawing aboard ship while locals left gold and retreated. Without meeting or speaking, both refrained from taking excess until a price was agreed. Later Arabic geographers described similar West African gold exchanges. Such dealings relied on shared rules, whose transmission suggests long-established trade.

\section{Through a Traveller's Eyes}
Read a short primary account and apply the tool.

In 1352, Ibn Battuta crossed the Sahara with a caravan. At Taghaza he described harsh conditions, salt-block houses and mosque, camel-skin roofs, and inhabitants growing no food, dependent on supplies from southern caravans, with salt their only capital.

In Mali, he described religious life, famously noting Qur'an memorisation enforced by fetters removed only after mastery. At a festival he saw children with chained feet and was told why. Subjects scattered dust on their heads before Mansa Sulayman, a court custom he openly disliked. He praised security allowing safe travel while criticising food and relations between women and men that violated his standards.

Ask our four questions. Who? A widely travelled Moroccan religious scholar, still an outsider. How much later? Essentially contemporaneous, giving weight. For whom? Arabic-world readers, hence home standards rewarding piety and condemning unfamiliar customs. \textbf{Praise and complaint both provide information}: praise documents security and education; complaint suggests actual encounter with difference rather than copying at home. Salt production, safety, ritual, and educational practice are credible observations; imperial origins and exhaustive regional products remain hearsay beyond his reach.

Add secondhand testimony. Cairo official and scholar al-Umari never visited West Africa but interviewed witnesses to Mansa Musa's 1324 pilgrimage through Cairo to Mecca: thousands of attendants and abundant charitable gold. He reports that spending and gifts depressed Cairo's gold price for years. One ruler's journey denting prices in a great Mediterranean commercial centre conveys West African gold's scale better than adjectives. Yet al-Umari is a relayer and figures may exaggerate: accept vast scale, not a particular number.

\section{Why Empires Rose and Departed}
Establish the facts. Ancient Ghana, centred in today's Mali and Mauritania, flourished roughly in the eighth to eleventh centuries; Mali in the thirteenth to sixteenth; Songhay in the fifteenth to sixteenth. Successive control of gold regions and trade changed capitals and ports. In 1591, Moroccan troops crossed the Sahara with firearms and defeated Songhay at Tondibi, ending the era of West African savanna empires.

Other African patterns coexisted. Aksum on the Ethiopian highlands minted coins and traded through the Red Sea in the first millennium. Christianity arrived around the fourth century, followed by over a thousand years of manuscript production. Eleven rock-hewn churches, traditionally attributed to King Lalibela in the twelfth and thirteenth centuries, remain living holy places. Swahili cities such as Kilwa, Mombasa, and Malindi exported inland gold and ivory towards Arabia, India, and China; Ibn Battuta called Kilwa among the best-built cities he had seen. Chinese ceramics still emerge from their ruins. Great Zimbabwe controlled inland gold regions linked to coastal outlets. Central kingdoms, loose city-state associations, and communities functioning without kings shared no single template and needed not resemble Europe.

Those are events. At least three nonequivalent explanations compete over why the savanna empires rose and fell.

\textbf{Explanation one: control of trade routes.} Whoever held gold-and-salt junctions taxed each camel-load and gold pouch, funded armies, and protected routes. Empires arose by holding nodes and changed when nodes changed hands. This explains regional succession neatly, not precise collapse dates. Profitable route control should not suddenly fail in 1591 unless something else changed.

\textbf{Explanation two: politics and war.} Human organisation creates empires and organisational failures destroy them. Later Mali suffered succession disputes and weak control of distant rulers; Songhay faced firearms that scattered its cavalry despite Morocco's small army. This supplies timing through crises and new weapons. Yet it produces chains of counterfactuals: what if Songhay had guns or a stronger successor? It sees events, not their underlying slow variables.

\textbf{Explanation three: climate and trade's slow shifts.} Sahelian farming and pastoralism are highly rainfall-sensitive, and centuries of drying strained finances. More importantly, fifteenth-century Portuguese coastal trade gave gold another exit south to sea rather than north across desert. Maritime routes bypassed savanna tollbooths. This broadly explains decline's window and the absence of later equally large savanna empires. Its limitation is that explaining everything through slow variables risks explaining nothing: cities sometimes flourished amid drought, and trans-Saharan salt trade survived Atlantic expansion into the early twentieth century.

Each holds a causal segment. Ecology and routes shape the slope; politics and war determine the collapse's moment. Distinguish imperial events, competing causal accounts, contemporaries' worldviews such as chroniclers blaming immoral rulers, and our judgements about exploitative taxation or firearms as progress. Mixing them turns history into propaganda.

\section[Further Challenge: History without Writing]{Further Challenge: Does Unwritten History Count?}
Our weightiest question comes from a revolution within history.

In the 1960s, Belgian-born Jan Vansina published the classic Oral Tradition, offending assumptions by systematically arguing that spoken traditions could be used as seriously as archives, provided they received equally rigorous examination.

His method reconstructs a transmission chain from witness through successive tellers to recorder. Who recorded it, after how many generations, each with which interests and roles? Then distinguish fixed and flexible elements. Griots carry structured king lists, genealogies, place names, and formulas, remarkably stable across centuries and regions, alongside redecorated dialogue, emotions, and scenes. Sundiata's framework can serve history while its embellishments serve literature.

Vansina's deeper impact challenged evidence itself. Why should paper automatically outrank minds? Paper does not misremember, but writers lie. Memory changes, but collectively guarded traditions have correction mechanisms: other griots can challenge a mistaken king list. Literate elites largely established writing's superior rank while monopolising writing. This does not discredit texts; it means \textbf{reliability should rest on method, not medium}.

Two comparisons show the force. Troy lived for over a millennium in Homeric oral tradition, often dismissed as myth, until nineteenth-century Schliemann followed epic clues to northwestern Asia Minor and excavated a city now generally linked to that remembered past. Europe's own roots were recovered from oral tradition. In South America, the Inca governed over ten million people with roads, treasury, and census but no general writing system. Quipu knots recorded numbers and affairs through colour, position, and form. Complex society plainly need not require writing.

Africa supplies another easily misjudged counterexample: sub-Saharan Africa was not without writing. Ethiopian Ge'ez appeared on inscriptions and coins around the fourth century, with continuous Christian manuscript traditions thereafter. Swahili cities wrote their language in Arabic script, preserving poetry and accounts. The familiar claim of an unwritten Africa has a false premise.

The challenge is double: weigh oral evidence methodically, and question who calibrated the scales.

\section{Today: Filling the Blank Page by Page}
This old question remains active nearby.

After Timbuktu's 2012 rescue, manuscripts are being registered, restored, and digitised for worldwide researchers. Accounts, contracts, and letters offer coveted social-history evidence. Their knowledge network becomes a searchable database rather than legend.

Consider museums. British troops captured Benin's capital in 1897, in present-day Nigeria, not today's country named Benin, and looted thousands of bronze and ivory works dispersed among European and American museums. From 2022, Germany led large-scale ownership returns to Nigeria, with others following under pressure. Returning objects also returns narrative authority: bronze reliefs record Benin court history, a cast history book taken by force and now reclaimed.

Maps offer a final caution. Most modern African boundaries were drawn by colonial rulers around the Berlin Conference, straight through peoples and around mineral regions, largely ignoring local historical geography. Do not project modern names backwards. Independent Ghana deliberately adopted ancient Ghana's name in 1957, though the empire lay over a thousand kilometres away in today's Mali and Mauritania: homage, not geographical continuity. Benin is similar. Nor can ancient Nigeria be one box containing Benin, Yoruba and Hausa cities, and communities without central kingship. Colonial lines framing precolonial worlds compound anachronism.

Beware the reverse extreme too. Renewed African history sometimes prompts online fantasies attributing every civilisation to Africa. This repeats denial's method: decide first, seek evidence later, merely reversing the conclusion. Corroboration, source scrutiny, and separating fact from judgement apply equally. Wronged histories need no compensatory myths; reality is grand enough without inflation.

Finally, classrooms still devote far less space to Africa than Europe in many countries, but change continues. Great Zimbabwe's stone walls appear on Zimbabwe's flag and money; Timbuktu libraries welcome visitors; young people of many backgrounds discover that their presumed precolonial blank continent is itself one of colonialism's most persistent legacies.

\section[History Carried by Memory Alone]{Your Question: If Memory Alone Carried History}
Return to the manuscript page.

It is a fortunate exception: written down and rescued, it survives to testify for a vast world. Much more---griots' king lists, market accounts, builders' names, caravan leaders' lifetime routes---left no paper, remaining in memory or disappearing.

Imagine historians three centuries hence studying your city after electronic records, paper, and photographs are destroyed. Only three things survive: elders recalling ancestors' stories, buried foundations and rubbish layers, and travel accounts by foreign visitors who stayed three months without knowing local dialects.

Can those historians write a credible account? What matters most among what they will miss? If someone declares that without written archives the city never existed, how would you refute that?

Go further: when we call history credible, do we trust paper or a method of tracing sources and cross-checking? Does storing history in memory make a society weaker, or simply use a format we have not learned to read?

There is no standard answer. But your reasons may use the same ruler as those who denied African history three centuries ago. Rather than imperial names you may forget, keep this habit: before using a ruler, ask who made it.
